\section{Introduction}
\label{submission:intro}

The modern deep learning paradigm is heavily anchored in pre-training followed by task-specific fine-tuning.
As a result, the community has seen an explosion of highly specialized expert models, each adapted from a shared pre-trained backbone to excel at a particular downstream task.
While these experts exhibit remarkable task-specific performance, hosting, serving, and maintaining a massive collection of individual models for each task is computationally and operationally prohibitive.
To address this, there has been a surge of interest in \emph{model merging}---the process of combining multiple independently fine-tuned expert models into a single unified multi-task network without additional training, backpropagation, or costly data curation \cite{wortsman2022robust, ilharco2022editing, yadav2023ties}.

At the heart of model merging is the concept of \emph{task vectors}, defined as the difference between a fine-tuned expert's weights and the pre-trained base model's weights \cite{ilharco2022editing}.
Under simple \emph{Task Arithmetic}, task vectors are linearly scaled and summed to produce a merged model.
Although elegant and direct, Task Arithmetic frequently suffers from severe performance degradation, or even complete functional collapse, when combining more than a few models.
This collapse is primarily driven by \emph{destructive task interference}: different expert models often push the same underlying parameter in opposite directions, resulting in sign conflicts that neutralize or distort the learned representations \cite{yadav2023ties}.

To mitigate task interference, recent state-of-the-art methods have introduced sophisticated, multi-stage heuristics.
The leading baseline, \mbox{TIES-Merging} \cite{yadav2023ties}, addresses sign conflicts by executing a four-step pipeline: (1) trimming the bottom $k\%$ of values by magnitude to remove redundant updates, (2) computing a consensus sign-vote across all models for each parameter index, (3) filtering out (zeroing) any updates that do not conform to the elected consensus sign, and (4) rescaling the remaining active values to preserve parameter energy.
Other methods, such as \mbox{DARE} \cite{yu2024dare}, employ randomized dropout and scaling, while test-time adaptive approaches like \mbox{SyMerge} \cite{yang2024symerge} optimize layer-wise merging coefficients on unlabeled test data using gradients.

\begin{figure}[t]
\centering
\includegraphics[width=\columnwidth]{AdaMerging.png}
\caption{Conceptual overview of model merging methods. (a) Task Arithmetic directly averages task vectors, leading to destructive interference when updates oppose each other. (b) \mbox{TIES-Merging} resolves conflicts using multi-stage heuristics: magnitude trimming, sign-voting consensus, non-conforming zeroing, and parameter rescaling. (c) Our proposed \textbf{\mbox{Winner-Take-All} Sign Election (WTA-Sign)} strips away this complexity. Under Occam's razor, update magnitude is a direct proxy for confidence. The expert with the largest absolute update at each parameter index simply elects the sign. Non-conforming updates are masked out, and conforming updates are averaged, with zero hyperparameters.}
\label{fig:concept}
\end{figure}

Under the guidance of \emph{Occam's razor}, we argue that the deep learning community has fallen into a pattern of needless complexity.
In this paper, we advocate for a return to simplicity.
We question whether the complex heuristics of \mbox{TIES-Merging}---specifically, the arbitrary trimming of weights, the multi-step sign voting consensus, and the hyperparameter-dependent rescaling---are truly necessary to resolve weight-space conflicts.
Instead, we propose that the magnitude of an expert's parameter update is a natural and sufficient proxy for its confidence.
When multiple experts conflict, the expert that has undergone the most significant update is highly likely to be the most critical for that specific weight's representational capacity.

Based on this minimalist philosophy, we introduce \textbf{\mbox{Winner-Take-All} Sign Election (WTA-Sign)}.
WTA-Sign resolves sign interference via an elegant, closed-form, training-free mechanism:
\begin{enumerate}
    \item \textbf{Winner Election:} At each individual parameter index, we identify the single expert model with the largest absolute update magnitude.
    \item \textbf{Sign Election:} The sign of this winning update is elected as the direction for the merged parameter.
    \item \textbf{Conformity Masking \& Averaging:} We construct a binary mask to filter out any updates from other experts that oppose this elected sign. We then average only the conforming updates.
\end{enumerate}

WTA-Sign is completely parameter-free and hyperparameter-free: it has no trimming thresholds to set, no random dropout probabilities to tune, and no rescaling multipliers to calibrate.
It requires no training, no test-time optimization, and is completely deterministic.
Yet, as our experiments demonstrate, it achieves superior performance over both Task Arithmetic and TIES-Merging.

We summarize our core contributions as follows:
\begin{itemize}
    \item \textbf{Minimalist Philosophy:} We challenge the trend of increasingly complex model merging heuristics, showing that excellent interference mitigation can be achieved by prioritizing simplicity and confidence-based magnitude.
    \item \textbf{The WTA-Sign Method:} We introduce WTA-Sign, a training-free, parameter-free, closed-form merging algorithm that resolves parameter conflicts by letting the most confident expert elect the sign of the merge.
    \item \textbf{Empirical Superiority:} We evaluate WTA-Sign against Model Soups, Task Arithmetic, and TIES-Merging on a multi-task setup consisting of \mbox{MNIST}, \mbox{SVHN}, and \mbox{CIFAR10} with a CLIP \mbox{ViT-B-32} backbone. WTA-Sign completely avoids the interference-driven collapse of Task Arithmetic, maintaining a top average accuracy of \textbf{14.19\%} (fully preserving zero-shot capabilities), and consistently outperforming TIES-Merging (12.92\%) across scaling sweeps.
    \item \textbf{Elegance of Implementation:} WTA-Sign is exceptionally simple to implement. We provide a fully parallelized, vectorized PyTorch implementation that spans only four lines of code and introduces zero computational or parameter overhead.
\end{itemize}
